% 原书页码：038；表格、表题与表注共用一个环境。
\begin{table}[!htbp]\phantomsection\label{table-57}
\centering\BookTableStyle
\caption*{表格1\quad 汤若望大寿贺文摘要}
\begin{tabularx}{\linewidth}{>{\centering\arraybackslash}p{24mm}|>{\hsize=.85\hsize\raggedright\arraybackslash}X|>{\hsize=1.15\hsize\raggedright\arraybackslash}X}
\hline
祝寿者 & 资历简介 & 贺文摘要* \\ \hline
金之俊\newline （1593—1676） & 明万历年间进士（1619年）；明末任礼部郎中。降清后，为太子太保、兵部尚书、左都御史、吏部尚书、国史院大学士、中和殿大学士。清决定为崇祯皇帝立碑，就由金之俊撰文。 & 贺文全长约800字。“如诸君子之言，先生殆以术而寓乎道者；余谓先生则以道而忘乎术者。盖先生之全乎道，非以术教而以身教者也。”“士大夫之朝夕习于先生者，钦其卑牧，饮其和醇。” \\ \hline
魏裔介\newline （1616—1686） & 1646年顺治朝进士；太子太傅、光禄大夫。官至太子太保、保和殿大学士。 & 贺文全长约1200字。“盖闻命世大才，经纶名教者，不必华宗夏士。”“以余观于道未汤老先生，殆器大神宏，而无愧于古之圣贤者欤。”“又发前人所未发，是以密合天行，特膺纶眷，所谓博物君子，学贯天人者，先生之谓也。”“谓先生为西海之儒，即中华之大儒可也。” \\ \hline
龚鼎孳\newline （1615—1673） & 1634年，18岁中进士。诗文大家，和钱谦益（1582—1664）、吴伟业（1609—1672）一起并称“江左三大家”。 & 贺文全长约1000字。在天文方面，把汤若望与僧一行和郭守敬并举。特别赞美了汤若望对顺治皇帝的劝谏。把汤与历史上的名臣相提并论。 \\ \hline
胡世安\newline （?—1663） & 1628年明代进士；武英殿大学士、太子太傅、兵部尚书。 & 贺文全长约400字。 \\ \hline
王崇简\newline （1602—1678） & 1643年明崇祯朝进士，翰林学士，任《明史》撰修。 & 贺文全长约600字。称汤若望到中国来“迪我中夏”。 \\ \hline
\end{tabularx}
\BookFloatNotes{* 引自《赠言》，1367—1375。}
\end{table}%
